% -----------------------------------------------------------------------------
% Document class
% scrartcl = clean article class for reports; bibliography in TOC, good tables
% -----------------------------------------------------------------------------
\documentclass[
  bibliography=totoc,
  captions=tableheading,
]{scrartcl}

% Fix minor KOMA compatibility issues
\usepackage{scrhack}

% Warn if recompilation needed (refs, TOC, etc.)
\usepackage[aux]{rerunfilecheck}

% Core math packages (align, symbols, extensions)
\usepackage{amsmath}
\usepackage{amssymb}
\usepackage{mathtools}

% Modern font handling (requires lualatex/xelatex)
\usepackage{fontspec}
\recalctypearea{} % recompute layout after font setup

% Language settings (hyphenation, captions)
\usepackage[english]{babel}

% Modern math fonts + ISO style (clean scientific notation)
\usepackage[
  math-style=ISO,
  bold-style=ISO,
  sans-style=italic,
  nabla=upright,
  partial=upright,
  mathrm=sym,
]{unicode-math}

\setmathfont{Latin Modern Math} % default math font
\setmathfont{XITS Math}[range={scr, bfscr}] % nicer script letters
\setmathfont{XITS Math}[range={cal, bfcal}, StylisticSet=1]

% Numbers & units formatting (\num, \SI, etc.)
\usepackage[
  locale=US,
  separate-uncertainty=true,
  per-mode=symbol-or-fraction,
  detect-all,
]{siunitx}

% Quotes (recommended with biblatex)
\usepackage[english=american]{csquotes}

% Float placement control (figures/tables positioning)
\usepackage{float}
\floatplacement{figure}{htbp}
\floatplacement{table}{htbp}

% Keep figures inside sections (prevents drifting)
\usepackage[section,below]{placeins}

% text wrapping around figures
\usepackage{wrapfig}      

% Caption styling (bold labels, smaller font)
\usepackage[
  labelfont=bf,
  font=small,
  width=0.9\textwidth,
]{caption}

% Subfigures (compare plots side-by-side)
\usepackage{subcaption}

% Include images
\usepackage{graphicx}

% Better tables (clean formatting + number alignment)
\usepackage{tabularray}
\UseTblrLibrary{booktabs, siunitx}

% Typography improvements (spacing, kerning)
\usepackage{microtype}

% Bibliography system (modern replacement for BibTeX)
\usepackage[backend=biber]{biblatex}
\addbibresource{../../../latex/references.bib}

% Clickable links + PDF metadata
\usepackage[
  unicode,
  pdfusetitle,
  pdfcreator={},
  pdfproducer={},
]{hyperref}

% Better PDF bookmarks
\usepackage{bookmark}

% Custom commands (needed for definitions below)
\usepackage{expl3}
\usepackage{xparse}

% Upright e (useful for exp, softmax, etc.)
\ExplSyntaxOn
\NewDocumentCommand \E {} {\symup{e}}
\ExplSyntaxOff

% Upright differential d (clean integrals)
\ExplSyntaxOn
\NewDocumentCommand \dif {m} {\mathinner{\symup{d} #1}}
\ExplSyntaxOff

% Vector command (bold vectors in math mode)
\ExplSyntaxOn
\let\vaccent=\v
\RenewDocumentCommand \v {}{
  \TextOrMath{\vaccent}{\symbf}
}
\ExplSyntaxOff

% Better spacing for scalable brackets
\usepackage{mleftright}

% Shortcuts for scalable brackets: \l( ... \r)
\ExplSyntaxOn
\let\ltext=\l
\RenewDocumentCommand \l {}{\TextOrMath{\ltext}{\mleft}}
\let\raccent=\r
\RenewDocumentCommand \r {}{\TextOrMath{\raccent}{\mright}}
\ExplSyntaxOff

% Absolute value and norm (auto-scaling with *)
\DeclarePairedDelimiter{\abs}{\lvert}{\rvert}
\DeclarePairedDelimiter{\norm}{\lVert}{\rVert}

% Definition equality symbols
\newcommand{\defeq}{\vcentcolon=}
\newcommand{\eqdef}{=\vcentcolon}

% No paragraph indent (clean report style)
\setlength{\parindent}{0pt}

\title{Task 08: Transformer Reconstruction of IceCube Interaction Positions}
\author{Akram Aki}
\date{\today}

\begin{document}

\maketitle

\section{Introduction}

The goal of this task was to reconstruct the two-dimensional neutrino interaction position from IceCube hit sequences.
Each event is represented by a variable-length list of measured Cherenkov-light hits.
For every hit, the available input features are the detection time and the detector-unit position,
\[
    (t, x, y),
\]
and the regression target is the true interaction position
\[
    (x_{\mathrm{pos}}, y_{\mathrm{pos}}).
\]
The main technical difficulty is that different events contain different numbers of hits.
Therefore, the implementation uses a custom batching function, padding, and a padding mask before applying a Transformer encoder.

\section{Method}

The dataset was loaded from the local IceCube 2D parquet files.
The train, validation, and test splits each contained \(10001\) events.
Although the split sizes are identical, the files are distinct; for example, the first target coordinate differs between the three splits.
The training events contained between \(1\) and \(68\) hits, with a median of \(19\) hits and a mean of \(19.89\) hits.
\autoref{fig:dataset-overview} shows the hit-count distribution, the one-dimensional target-coordinate distributions, and the two-dimensional target density.

\begin{figure}[H]
    \centering
    \includegraphics[width=0.98\textwidth]{../build/dataset_overview.pdf}
    \caption{Overview of the IceCube data used in Task 08. The upper-left panel shows the number of hits per event for the three splits. The remaining panels show the training target distributions in \(x\), in \(y\), and in the two-dimensional interaction plane.}
    \label{fig:dataset-overview}
\end{figure}
\FloatBarrier

All input features and target coordinates were standardized using statistics from the training split only.
The same constants were then applied to the validation and test splits.
This avoids leaking information from validation or test data into the training procedure.
The model was trained in normalized coordinate space, while final metrics were computed after transforming predictions and labels back to meters.

The custom collate function converts each event from the Awkward-array layout \([3, N_{\mathrm{hits}}]\) to the PyTorch sequence layout \([N_{\mathrm{hits}}, 3]\).
Within one batch, all hit tensors are concatenated and the original sequence lengths are stored.
Inside the model, these lengths are used to split the hits back into events, pad the events to the longest sequence in the batch, and construct a boolean padding mask.
The mask is passed to \texttt{nn.TransformerEncoder} so that padded positions do not contribute to attention.
After the encoder, masked mean pooling averages only over real hits and produces one event-level vector.

The Transformer architecture is summarized in \autoref{tab:model-config}.
It first projects the three hit features to a hidden dimension of \(64\), applies two Transformer encoder layers with four attention heads, and maps the pooled event representation to the two target coordinates.
The model has \(71618\) trainable parameters.

\begin{table}[H]
    \centering
    \begin{tabular}{ll}
        \toprule
        Quantity & Value \\
        \midrule
        Input features per hit & \(t, x, y\) \\
        Targets & \(x_{\mathrm{pos}}, y_{\mathrm{pos}}\) \\
        Batch size & 64 \\
        Maximum epochs & 50 \\
        Early-stopping patience & 6 \\
        Hidden dimension & 64 \\
        Attention heads & 4 \\
        Transformer encoder layers & 2 \\
        Feed-forward dimension & 128 \\
        Dropout & 0.05 \\
        Optimizer & AdamW \\
        Learning rate & \(10^{-3}\) \\
        Weight decay & \(10^{-4}\) \\
        Loss & Mean squared error \\
        Trainable parameters & 71618 \\
        \bottomrule
    \end{tabular}
    \caption{Transformer and training configuration used for the Task 08 notebook run.}
    \label{tab:model-config}
\end{table}

\section{Results}

\autoref{fig:loss-curves} shows the training and validation losses in normalized coordinate space.
Training was configured for up to \(50\) epochs, but early stopping stopped the run after epoch \(35\).
The best validation checkpoint was selected at epoch \(29\), where the validation MSE reached \(0.0805\).
The final test MSE in normalized space was \(0.0774\).

\begin{figure}[H]
    \centering
    \includegraphics[width=0.82\textwidth]{../build/transformer_loss_curves.pdf}
    \caption{Training and validation loss curves for the Transformer encoder. The vertical line marks the best validation checkpoint, which was used for the final test evaluation.}
    \label{fig:loss-curves}
\end{figure}
\FloatBarrier

After denormalizing the predictions and labels, the model reached the metrics shown in \autoref{tab:test-metrics}.
The mean two-dimensional position error was \(\SI{1.3492}{m}\), and the median error was \(\SI{0.8767}{m}\).
The \(R^2\) score of \(0.9233\) indicates that the model explains most of the coordinate variance in the held-out test set.

\begin{table}[H]
    \centering
    \begin{tabular}{ll}
        \toprule
        Metric & Value \\
        \midrule
        Test MSE, normalized space & 0.0774 \\
        Physical-space MSE & 1.9051 \\
        MAE in \(x\) & \(\SI{0.7802}{m}\) \\
        MAE in \(y\) & \(\SI{0.9324}{m}\) \\
        Mean position error & \(\SI{1.3492}{m}\) \\
        Median position error & \(\SI{0.8767}{m}\) \\
        \(R^2\) score & 0.9233 \\
        \bottomrule
    \end{tabular}
    \caption{Final test-set metrics after converting the normalized model outputs back to meters.}
    \label{tab:test-metrics}
\end{table}

\autoref{fig:pred-true} compares predicted and true coordinates separately.
Most points lie close to the perfect-prediction diagonal, but the spread is slightly larger for the \(y\) coordinate, consistent with the larger \(y\)-MAE in \autoref{tab:test-metrics}.

\begin{figure}[H]
    \centering
    \includegraphics[width=0.95\textwidth]{../build/predicted_vs_true_positions.pdf}
    \caption{Predicted versus true interaction coordinates on the test set. The diagonal line indicates perfect reconstruction.}
    \label{fig:pred-true}
\end{figure}
\FloatBarrier

\autoref{fig:error-summary} shows the full distribution of Euclidean position errors.
The distribution is strongly concentrated below a few meters, but it has a tail of harder events with larger reconstruction errors.

\begin{figure}[H]
    \centering
    \includegraphics[width=0.78\textwidth]{../build/position_error_histogram.pdf}
    \caption{Distribution of two-dimensional position reconstruction errors on the test set.}
    \label{fig:error-summary}
\end{figure}
\FloatBarrier

Finally, \autoref{fig:error-hit-count} checks whether events with more detected hits are easier to reconstruct.
The binned trend decreases with increasing hit count, which is physically plausible because more hits provide more timing and geometry information to the Transformer.

\begin{figure}[H]
    \centering
    \includegraphics[width=0.78\textwidth]{../build/position_error_vs_hit_count.pdf}
    \caption{Position error as a function of the number of hits in the test event. The red markers show binned means with standard-deviation bars.}
    \label{fig:error-hit-count}
\end{figure}
\FloatBarrier

\section{Conclusion}

The Transformer encoder successfully handled the variable-length IceCube hit sequences using padding masks and masked mean pooling.
The model achieved a mean test position error of \(\SI{1.3492}{m}\) and a median error below one meter.
The error-versus-hit-count diagnostic also shows that events with more detected photons are generally reconstructed more accurately.

Possible improvements include tuning the model size, adding positional or time-order embeddings, using a learning-rate schedule, or comparing masked mean pooling with a learned regression token.
One important limitation of this run is the relatively small local training split of \(10001\) events.
The dataset description linked in the task material suggests that a much larger IceCube sample is available, so training on more events would be a natural next step and could improve generalization.
It would also be useful to compare the Transformer directly against the previous graph-neural-network reconstruction model on the same split.
All code used for data loading, training, evaluation, and plotting is available in the accompanying
\href{https://github.com/AkramAki/Advanced_DeepLearning_Seminar}{GitHub repository}.

\printbibliography

\end{document}
